\documentclass{article}
\usepackage[utf8]{inputenc}
\usepackage{geometry}
\usepackage{enumitem}
\usepackage{hyperref}
\usepackage{xcolor}
\usepackage{graphicx}
\usepackage{array}
\usepackage{multirow}
\usepackage{amsmath, amssymb}
\usepackage{chemfig}
\usepackage{mhchem}
\usepackage{tikz}
\usepackage{setspace}
\usepackage{soul}
\usepackage{mathtools}
\usepackage{booktabs}
\usepackage{chemformula}
\usepackage{siunitx}
\usetikzlibrary{positioning, calc}
\usepackage{longtable}
\geometry{margin=1in}
\setlength{\parskip}{0.5em}
\usepackage{cancel}

\begin{document}

\begin{center}
    {\LARGE\textbf{SI Session Planning Form}}
\end{center}

\
\noindent
\begin{flushleft}
\textbf{SI Leader:} Brandon Cobb\\
\textbf{Session Week:} 10\\
\textbf{Session\#:} 10\\
\textbf{Instructor:} Professor Tramontozzi\\
\textbf{Old Plans:}\href{https://macombedu.sharepoint.com/:b:/s/ASC-SupplementalInstructionEmbeddedTutoring/EeSYqCwT_KhMolgJvXbVEjgBOyF_DevvS8BgCm1zPkQQJg?e=QJTt4o}{SharePoint}\\
\textbf{Session Goal:} Students will review and apply concepts of stoichiometry, including balancing chemical equations, calculating moles, and using mole ratios to relate reactants and products. They will apply gas laws (Boyle’s, Charles’s, Avogadro’s) and the kinetic molecular theory to understand relationships among pressure, volume, temperature, and number of gas particles. Additionally, students will analyze solutions, calculating concentrations (mass percent, volume percent, molarity), applying dilution concepts, and using solubility rules to predict saturation and precipitation.
\end{flushleft}

\vspace{1em}
\section*{\underline{Opener}}
{\centering\large\textbf{\hl{Take Attendance}}\par}\vspace{-0.25\baselineskip}
\noindent
\begin{tabular}{|p{4.2cm}|p{9.8cm}|}
\hline
\textbf{Collaborative Learning Techniques} & Group Discussion\\
\hline
\textbf{Learning Strategy} & Lecture Review\\
\hline
\textbf{Time} & 12:00---12:10\\
\hline
\textbf{Materials} & Notes, homework, whiteboard, markers, ready to speak, periodic table \\
\hline
\textbf{Lecture/Textbook/Old Plans/Slide References} & Class roster, Chapter 7: Gases, Liquids and Solids, Chapter 8: Solutions  \\
\hline
\end{tabular}


\vspace{0.5cm}
{\large\textbf{Definitions}}
\\
\begin{tabular}{|p{0.9\linewidth}|}
\hline
\begin{minipage}[t]{0.9\linewidth}
\begin{enumerate}
    \item\textbf{Group Discussion:} One person in the group is asked to present on a topic or review material for the group and then lead the discussion for the group. This person should not be the regular group leader.
   \item\textbf{Lecture Review:} As a group summarize the lecture from the previous class. You may have to provide prompts for the students. For example, "The first concept discussed was Civil Liberties and Public Policy, what did the profressor highlight regarding this?" You may want to ask them to try summarizing without looking at their notes; however, if they are having a difficult time remembering, tell them to refer to their notes.
\end{enumerate}
\end{minipage} \\
\hline
\end{tabular}
\newpage
\noindent
{\large\textbf{Instructions (please provide a\hl{DETAILED} activity breakdown below (and/or attached}}\\

\begin{tabular}{|p{0.9\linewidth}|}
\hline
\vspace{0.2cm}
\textbf{Lecture Review:}
\begin{enumerate}
\item During te first 10-15 minutes of the SI session. have the students summarize the most recent lecture, or have them identify the key words from that lecture.
\item Give students three minutes to find support in their lecture notes for a given generalization.
\item Have the students predict the direction of future lectures based upon the past lectures.
\item Have the student arrange terms from lecture and text into a structured outline.
\item Reinforce new terms or important information by using clearly constructed handouts (can be complete or nearly complete at the beginning of the term but should gradually require more and more filling in as the group becomes more accustomed to working together.
\item Review material from previous sessions and lectures.
\item Take a couple of minutes at the end of the SI session to summarize the main idea convered during a session. Asl the students to help summarize.
\item Have students write a one paragraph summary of the lecture. List the new vocabulary terms introduced with this lecture.
\item Formulate potential exam questions based on the main ideas from the lecture.
\item Formulate pontential answers from details in the lecture notes.
\end{enumerate} \\
\hline
\end{tabular}

\newpage
\noindent
{\large\textbf{Checking for Understanding (questions \& answers)} \\
\vspace{0.2cm}
\begin{tabular}{|p{0.9\linewidth}|}
\hline
\begin{itemize}
\item How do you convert grams of a substance to moles using its molar mass?
\item How do mole ratios in a balanced equation allow prediction of product formation?
\item Why is it important to identify excess and limiting reactants before calculating yields?
\item If a reaction produces less product than predicted, what might be the reason?
\item How can you determine how much of a reactant is required to completely react with another?
\item How do you calculate total pressure in a container with multiple gases?
\item How do the mole fractions of individual gases affect their partial pressures?
\item How does opening a connection between two gas containers affect the pressures in each container?
\end{itemize} \\
\hline
\end{tabular}

\newpage
\noindent
{\large\textbf{Checking for Understanding (questions \& answers) continued...} \\
\vspace{0.2cm}
\begin{tabular}{|p{0.9\linewidth}|}
\hline
\begin{itemize}
\item How does compressing a gas (reducing its volume) affect its pressure at constant temperature?
\item If the pressure of a gas is halved at constant temperature, what happens to its volume?
\item Why does a diver’s bubble expand as it rises to the surface?
\item How does increasing the number of moles of gas in a container at constant temperature and pressure affect its volume?
\item How does decreasing moles of gas affect volume in a sealed balloon?
\item Why does Avogadro’s law allow prediction of gas volume based on moles?
\item How does increasing the number of gas moles in a rigid container change the pressure?
\item How does removing gas molecules affect pressure in a fixed-volume container?
\item Why is the pressure proportional to the number of particles in a container?
\item How does heating a gas at constant pressure affect its volume?
\item How does cooling a gas at constant pressure affect its volume?
\item How does the Kelvin temperature scale relate to gas expansion and contraction?
\item How do you calculate mass percent of a solute in a solution?
\item How do you calculate volume percent of a solute in a solution?
\item How is molarity defined and calculated?
\item How does dilution affect the molarity of a solution?
\item Why is it important to consider solubility rules when predicting whether a precipitate will form?
\item How do you determine if a solution is saturated, unsaturated, or supersaturated?
\end{itemize} \\
\hline
\end{tabular}
\newpage

\section*{\underline{Activity 1}}
\vspace{0.2cm}
\begin{tabular}{|p{4.2cm}|p{9.8cm}|}
\hline
\textbf{Collaborative Learning Techniques} & Jigsaw\\
\hline
\textbf{Learning Strategy} & Predict Test Questions \\
\hline
\textbf{Time} & 12:10---12:30\\
\hline
\textbf{Materials} & Pencil, eraser, calculator, periodic table \\
\hline
\textbf{Lecture/Textbook/Old Plans/Slide References} & Chapter 7: Gases, Liquids and Solids, Chapter 8: Solutions \\
\hline
\end{tabular}
\vspace{0.5cm}

{\large\textbf{Definition}}\\
\noindent
\begin{tabular}{|p{0.9\linewidth}|}
\hline
\begin{enumerate}
    \item\textbf{Jigsaw:} Jigsaws, when used properly, make the group as a whole dependent upon all of them in subgroups. Each group provides a peice of the puzzle. Group members are broken into smaller groups. Each small group works on some aspect of the same problem, question, or issue. They then share their part of the puzzle with the large group.
    \item\textbf{Predict Test Questions:}  Put students in groups of 2 or 3 and assign them to write a test question for a specific topic, ensuring that all topics have been convered. Ask students to write their questions on the board or on an overhead for discussion (would the professor ask the question? What is the answer? etc.). Students will have the benefit of learning to think like the teacher and they'll be able to see additional questions that other students have writen.
 
\end{enumerate}\\
\hline
\end{tabular}
\vspace{0.2cm}
\newpage
{\large\textbf{Instructions (please provide a\hl{DETAILED} activity breakdown below (and/or attached}}\\
\noindent
\begin{tabular}{|p{0.9\linewidth}|}
\hline
\vspace{0.2cm}
\begin{enumerate}
    \item \textbf{Distribute Worksheets:} Each student receives a worksheet containing chemistry questions. The problems cover a range of difficulty so that students can decide which ones best match their needs.
    
    \item \textbf{Student Choice of Starting Point:} Students look over the worksheet and select their own starting problem rather than being assigned a fixed order. This allows them to challenge themselves where they feel it is most beneficial.
    
    \item \textbf{Staggered Timing:} Students work in pairs. Each partner begins at a different point (for example, one may start at Problem 1 while the other starts at Problem 3). They continue by alternating every three questions. This staggered structure ensures that students will be asking for help or requesting solutions at different times, keeping the class flow balanced.
    
    \item \textbf{Independent Work:} Students attempt their chosen problems individually, writing down both reasoning and answers. Emphasis is placed on documenting thought process as well as the result.
    
    \item \textbf{Help and Solution Reveal:} After the set time for each problem expires, the SI leader either walks through the solution on the whiteboard or directs students to check the answer key. Students are encouraged to ask clarifying questions at the timed checkpoints rather than all at once.
    
    \item \textbf{Partner Collaboration:} Once a solution is reviewed, partners share their reasoning for their respective problems. They ask each other questions, offer corrections, and compare how their approaches aligned with the solution.
    
    \item \textbf{Rotation of Problems:} After every three questions, students switch to new problems, continuing the staggered order so that they stay offset from each other. This ensures that each student engages with a variety of problems across the worksheet.
    
    \item \textbf{Class-Wide Debrief:} The SI leader pauses periodically to highlight frequently asked questions, difficult concepts, or common errors observed during the staggered work. Whole-class solutions are discussed when many students identify the same challenge.
    
    \item \textbf{Reflection:} At the close, students note which problems gave them the most insight, what strategies helped, and what topics they still want clarified. This reflection helps focus future review sessions.
\end{enumerate}\\
\hline
\end{tabular}
\newpage

\noindent
{\large\textbf{Content (fully worked out problems \& answers) continued...}}\\[0.2em]
\vspace{0.2cm}

\vspace{0.2cm}
\noindent\begin{tabular}{|p{0.9\linewidth}|}
\hline
\textbf{(1)} Which of the following solutes is generally soluble in water according to solubility rules?\\[0.2cm]
\begin{minipage}[t]{0.85\linewidth}
\begin{enumerate}[label=\alph*.]
\item BaSO$_4$
\item AgCl
\item KNO$_3$
\item PbI$_2$
\end{enumerate}
\end{minipage}\\[0.2cm]
\hline
\end{tabular}

\noindent\begin{tabular}{|p{0.9\linewidth}|}
\hline
\textbf{(2)} A chemist has 50.0~mL of 2.0~M HCl solution. She adds water to dilute it to a final volume of 200.0~mL. What is the molarity of the diluted solution?\\[0.2cm]
\begin{minipage}[t]{0.85\linewidth}
\begin{enumerate}[label=\alph*.]
\item 0.25~M
\item 0.50~M
\item 1.0~M
\item 2.0~M
\end{enumerate}
\end{minipage}\\[0.2cm]
\hline
\end{tabular}

\noindent\begin{tabular}{|p{0.9\linewidth}|}
\hline
\textbf{(3)} A balloon has a volume of 2.0~L at 300~K. If it is heated to 600~K at constant pressure, what will the new volume be?\\[0.2cm]
\begin{minipage}[t]{0.85\linewidth}
\begin{enumerate}[label=\alph*.]
\item 1.0~L
\item 2.0~L
\item 3.0~L
\item 4.0~L
\end{enumerate}
\end{minipage}\\[0.2cm]
\hline
\end{tabular}

\vspace{0.2cm}
\noindent\begin{tabular}{|p{0.9\linewidth}|}
\hline
\textbf{(4)} A diver's breathing mixture contains helium and oxygen at a total pressure of 8.0~atm. If the partial pressure of oxygen is 1.6~atm, what fraction of the total gas mixture is oxygen?\\[0.2cm]
\begin{minipage}[t]{0.85\linewidth}
\begin{enumerate}[label=\alph*.]
\item 0.10
\item 0.20
\item 0.50
\item 0.80
\end{enumerate}
\end{minipage}\\[0.2cm]
\hline
\end{tabular}
\newpage

\vspace{0.2cm}
\noindent\begin{tabular}{|p{0.9\linewidth}|}
\hline
\textbf{(5)} A 5.0~L sample of hydrogen gas is collected at 310~K. If it is cooled to 155~K at constant pressure, what is the new volume?\\[0.2cm]
\begin{minipage}[t]{0.85\linewidth}
\begin{enumerate}[label=\alph*.]
\item 1.0~L
\item 2.5~L
\item 5.0~L
\item 10.0~L
\end{enumerate}
\end{minipage}\\[0.2cm]
\hline
\end{tabular}

\noindent\begin{tabular}{|p{0.9\linewidth}|}
\hline
\textbf{(6)} A solvent is the substance present in the largest amount in a solution.\\[0.2cm]
\textbf{True} \hspace{1cm} \textbf{False}\\[0.2cm]
\hline
\end{tabular}

\vspace{0.2cm}
\noindent\begin{tabular}{|p{0.9\linewidth}|}
\hline
\textbf{(7)} A sample of gas has a volume of 3.0~L at 27~°C. What will its volume be at 127~°C if pressure remains constant?\\[0.2cm]
\begin{minipage}[t]{0.85\linewidth}
\begin{enumerate}[label=\alph*.]
\item 1.5~L
\item 3.0~L
\item 4.0~L
\item 5.0~L
\end{enumerate}
\end{minipage}\\[0.2cm]
\hline
\end{tabular}

\vspace{0.2cm}
\noindent\begin{tabular}{|p{0.9\linewidth}|}
\hline
\textbf{(8)} A mixture of helium and argon gases in a sealed 10~L container has a total pressure of 5.0~atm. If the partial pressure of helium is 2.0~atm, what is the partial pressure of argon?\\[0.2cm]
\begin{minipage}[t]{0.85\linewidth}
\begin{enumerate}[label=\alph*.]
\item 3.0~atm
\item 7.0~atm
\item 2.5~atm
\item 5.0~atm
\end{enumerate}
\end{minipage}\\[0.2cm]
\hline
\end{tabular}
\newpage

\vspace{0.2cm}
\noindent\begin{tabular}{|p{0.9\linewidth}|}
\hline
\textbf{(9)} A gas sample occupies 400~mL at 273~K. What volume will it occupy at 546~K if pressure remains constant?\\[0.2cm]
\begin{minipage}[t]{0.85\linewidth}
\begin{enumerate}[label=\alph*.]
\item 100~mL
\item 200~mL
\item 400~mL
\item 800~mL
\end{enumerate}
\end{minipage}\\[0.2cm]
\hline
\end{tabular}

\noindent\begin{tabular}{|p{0.9\linewidth}|}
\hline
\textbf{(10)} A drop of water beads up on a freshly waxed car surface. Which statement best explains why this happens?\\[0.2cm]
\begin{minipage}[t]{0.85\linewidth}
\begin{enumerate}[label=\alph*.]
\item Water's adhesion to the wax is stronger than its cohesion.
\item Hydrogen bonding causes strong cohesion within water.
\item London forces between wax molecules attract water.
\item Dipole-induced dipole forces repel water from wax.
\end{enumerate}
\end{minipage}\\[0.2cm]
\hline
\end{tabular}

\vspace{0.2cm}
\noindent\begin{tabular}{|p{0.9\linewidth}|}
\hline
\textbf{(11)} Carlos opens a bottle of perfume and quickly smells it from across the room. Which intermolecular force primarily determines how easily the perfume molecules enter the air?\\[0.2cm]
\begin{minipage}[t]{0.85\linewidth}
\begin{enumerate}[label=\alph*.]
\item Hydrogen bonding
\item Ionic bonding
\item London dispersion forces
\item Metallic bonding
\end{enumerate}
\end{minipage}\\[0.2cm]
\hline
\end{tabular}

\vspace{0.2cm}
\noindent\begin{tabular}{|p{0.9\linewidth}|}
\hline
\textbf{(12)} A student notices that hexane and water do not mix when shaken together. Which type of intermolecular forces explains this behavior?\\[0.2cm]
\begin{minipage}[t]{0.85\linewidth}
\begin{enumerate}[label=\alph*.]
\item London dispersion forces dominate in hexane, but water is polar.
\item Both substances form hydrogen bonds with each other.
\item Hexane is polar, and water is nonpolar.
\item Dipole-dipole interactions attract the two liquids together.
\end{enumerate}
\end{minipage}\\[0.2cm]
\hline
\end{tabular}
\newpage

\vspace{0.2cm}
\noindent\begin{tabular}{|p{0.9\linewidth}|}
\hline
\textbf{(13)} A saturated solution of NaNO$_3$ contains 90~g solute in 100~g water at 50~°C. If 10~g more NaNO$_3$ is added, what happens?\\[0.2cm]
\begin{minipage}[t]{0.85\linewidth}
\begin{enumerate}[label=\alph*.]
\item Dissolves completely
\item Remains undissolved
\item Reacts with water
\item Turns into a supersaturated solution spontaneously
\end{enumerate}
\end{minipage}\\[0.2cm]
\hline
\end{tabular}

\vspace{0.2cm}
\noindent\begin{tabular}{|p{0.9\linewidth}|}
\hline
\textbf{(14)} A weather balloon at 1.20~atm has a volume of 100~L near the ground. When it rises to a region where the pressure is 0.900~atm, what is its new volume (temperature constant)?\\[0.2cm]
\begin{minipage}[t]{0.85\linewidth}
\begin{enumerate}[label=\alph*.]
\item 75.0~L
\item 90.0~L
\item 100.~L
\item 133~L
\end{enumerate}
\end{minipage}\\[0.2cm]
\hline
\end{tabular}

\vspace{0.2cm}
\noindent\begin{tabular}{|p{0.9\linewidth}|}
\hline
\textbf{(15)} A scuba diver breathes from a tank where the air occupies 12.0~L at 200~atm. If the diver uses half of the air (6.00~L) while underwater at the same temperature, what is the new pressure in the tank?\\[0.2cm]
\begin{minipage}[t]{0.85\linewidth}
\begin{enumerate}[label=\alph*.]
\item 50~atm
\item 100~atm
\item 200~atm
\item 400~atm
\end{enumerate}
\end{minipage}\\[0.2cm]
\hline
\end{tabular}

\vspace{0.2cm}
\noindent\begin{tabular}{|p{0.9\linewidth}|}
\hline
\textbf{(16)} In a sugar-water mixture, sugar is the solvent and water is the solute.\\[0.2cm]
\textbf{True} \hspace{1cm} \textbf{False}\\[0.2cm]
\hline
\end{tabular}
\newpage

\noindent\begin{tabular}{|p{0.9\linewidth}|}
\hline
\textbf{(17)} A diver's air bubble has a volume of 3.0~L at a depth where the pressure is 2.0~atm. If the bubble rises to the surface where the pressure is 1.0~atm, what will its new volume be (assuming constant temperature)?\\[0.2cm]
\begin{minipage}[t]{0.85\linewidth}
\begin{enumerate}[label=\alph*.]
\item 1.5~L
\item 2.0~L
\item 3.0~L
\item 6.0~L
\end{enumerate}
\end{minipage}\\[0.2cm]
\hline
\end{tabular}

\vspace{0.2cm}
\noindent\begin{tabular}{|p{0.9\linewidth}|}
\hline
\textbf{(18)} In percent by volume concentration, both solute and solvent volumes are measured in the same units.\\[0.2cm]
\textbf{True} \hspace{1cm} \textbf{False}\\[0.2cm]
\hline
\end{tabular}

\vspace{0.2cm}
\noindent\begin{tabular}{|p{0.9\linewidth}|}
\hline
\textbf{(19)} When comparing HF and HCl, HF has a much higher boiling point. Which force best explains this difference?\\[0.2cm]
\begin{minipage}[t]{0.85\linewidth}
\begin{enumerate}[label=\alph*.]
\item Dipole-dipole interactions
\item Hydrogen bonding
\item London dispersion forces
\item Covalent bonding
\end{enumerate}
\end{minipage}\\[0.2cm]
\hline
\end{tabular}

\noindent\begin{tabular}{|p{0.9\linewidth}|}
\hline
\textbf{(20)} A piston traps 4.0~L of gas at 1.5~atm. The piston is pushed so that the gas pressure increases to 3.0~atm. What is the new volume of the gas?\\[0.2cm]
\begin{minipage}[t]{0.85\linewidth}
\begin{enumerate}[label=\alph*.]
\item 1.0~L
\item 2.0~L
\item 3.0~L
\item 6.0~L
\end{enumerate}
\end{minipage}\\[0.2cm]
\hline
\end{tabular}
\newpage

\vspace{0.2cm}
\noindent\begin{tabular}{|p{0.9\linewidth}|}
\hline
\textbf{(21)} The concentration of a solution in percent by mass is calculated using the mass of solute divided by the total mass of the solution.\\[0.2cm]
\textbf{True} \hspace{1cm} \textbf{False}\\[0.2cm]
\hline
\end{tabular}

\vspace{0.2cm}
\noindent\begin{tabular}{|p{0.9\linewidth}|}
\hline
\textbf{(22)} A steel tank of fixed volume contains 1.0~mol of nitrogen gas at 2.0~atm. If another 1.0~mol of nitrogen is added at the same temperature, what is the new pressure in the tank?\\[0.2cm]
\begin{minipage}[t]{0.85\linewidth}
\begin{enumerate}[label=\alph*.]
\item 1.0~atm
\item 2.0~atm
\item 3.0~atm
\item 4.0~atm
\end{enumerate}
\end{minipage}\\[0.2cm]
\hline
\end{tabular}

\vspace{0.2cm}
\noindent\begin{tabular}{|p{0.9\linewidth}|}
\hline
\textbf{(23)} A sealed container holds 5.00~L of oxygen gas at 1.00~mol. If half the gas leaks out while temperature and pressure stay constant, what will the new volume be?\\[0.2cm]
\begin{minipage}[t]{0.85\linewidth}
\begin{enumerate}[label=\alph*.]
\item 1.25~L
\item 2.50~L
\item 5.00~L
\item 10.0~L
\end{enumerate}
\end{minipage}\\[0.2cm]
\hline
\end{tabular}

\vspace{0.2cm}
\noindent\begin{tabular}{|p{0.9\linewidth}|}
\hline
\textbf{(24)} What is the molarity of a solution containing 5.00~mol of KBr in 2.00~L of solution?\\[0.2cm]
\begin{minipage}[t]{0.85\linewidth}
\begin{enumerate}[label=\alph*.]
\item 2.00~M
\item 2.50~M
\item 5.00~M
\item 10.0~M
\end{enumerate}
\end{minipage}\\[0.2cm]
\hline
\end{tabular}
\newpage

\vspace{0.2cm}
\noindent\begin{tabular}{|p{0.9\linewidth}|}
\hline
\textbf{(25)} Molarity expresses concentration as moles of solute per liter of solvent.\\[0.2cm]
\textbf{True} \hspace{1cm} \textbf{False}\\[0.2cm]
\hline
\end{tabular}

\vspace{0.2cm}
\noindent\begin{tabular}{|p{0.9\linewidth}|}
\hline
\textbf{(26)} An unsaturated solution contains less solute than the maximum that can dissolve at that temperature.\\[0.2cm]
\textbf{True} \hspace{1cm} \textbf{False}\\[0.2cm]
\hline
\end{tabular}

\noindent\begin{tabular}{|p{0.9\linewidth}|}
\hline
\textbf{(27)} A scuba tank contains 0.80~mol of air in a 6.0~L tank. How many moles of air would fill a 12.0~L tank at the same temperature and pressure?\\[0.2cm]
\begin{minipage}[t]{0.85\linewidth}
\begin{enumerate}[label=\alph*.]
\item 0.40~mol
\item 0.80~mol
\item 1.6~mol
\item 3.2~mol
\end{enumerate}
\end{minipage}\\[0.2cm]
\hline
\end{tabular}

\noindent\begin{tabular}{|p{0.9\linewidth}|}
\hline
\textbf{(28)} A scuba tank contains oxygen at 3.0~atm and nitrogen at 6.0~atm. What is the total pressure in the tank according to Dalton's Law?\\[0.2cm]
\begin{minipage}[t]{0.85\linewidth}
\begin{enumerate}[label=\alph*.]
\item 2.0~atm
\item 3.0~atm
\item 6.0~atm
\item 9.0~atm
\end{enumerate}
\end{minipage}\\[0.2cm]
\hline
\end{tabular}
\newpage

\vspace{0.2cm}
\noindent\begin{tabular}{|p{0.9\linewidth}|}
\hline
\textbf{(29)} According to solubility rules, most silver chloride (AgCl) salts are soluble in water.\\[0.2cm]
\textbf{True} \hspace{1cm} \textbf{False}\\[0.2cm]
\hline
\end{tabular}

\vspace{0.2cm}
\noindent\begin{tabular}{|p{0.9\linewidth}|}
\hline
\textbf{(30)} A gas sample in a steel container has a volume of 12.0~L at 600.~mmHg. If the pressure is reduced to 400.~mmHg, what is the new volume (temperature constant)?\\[0.2cm]
\begin{minipage}[t]{0.85\linewidth}
\begin{enumerate}[label=\alph*.]
\item 8.0~L
\item 12.0~L
\item 18.0~L
\item 20.0~L
\end{enumerate}
\end{minipage}\\[0.2cm]
\hline
\end{tabular}

\vspace{0.2cm}
\noindent\begin{tabular}{|p{0.9\linewidth}|}
\hline
\textbf{(31)} A solution with visible solid particles remaining after stirring is considered supersaturated.\\[0.2cm]
\textbf{True} \hspace{1cm} \textbf{False}\\[0.2cm]
\hline
\end{tabular}

\vspace{0.2cm}
\noindent\begin{tabular}{|p{0.9\linewidth}|}
\hline
\textbf{(32)} A saturated solution can dissolve more solute if the temperature is increased.\\[0.2cm]
\textbf{True} \hspace{1cm} \textbf{False}\\[0.2cm]
\hline
\end{tabular}
\newpage

\vspace{0.2cm}
\noindent\begin{tabular}{|p{0.9\linewidth}|}
\hline
\textbf{(33)} A chemist mixes CO$_2$ and N$_2$ gases in a flask. If CO$_2$ exerts a pressure of 0.750~atm and N$_2$ exerts 0.950~atm, what is the total pressure?\\[0.2cm]
\begin{minipage}[t]{0.85\linewidth}
\begin{enumerate}[label=\alph*.]
\item 0.200~atm
\item 1.70~atm
\item 0.850~atm
\item 1.00~atm
\end{enumerate}
\end{minipage}\\[0.2cm]
\hline
\end{tabular}

\vspace{0.2cm}
\noindent\begin{tabular}{|p{0.9\linewidth}|}
\hline
\textbf{(34)} A balloon inflated to 10.0~L at 293~K is placed in a freezer at 263~K. What is the new volume if pressure is constant?\\[0.2cm]
\begin{minipage}[t]{0.85\linewidth}
\begin{enumerate}[label=\alph*.]
\item 8.00~L
\item 9.00~L
\item 10.0~L
\item 11.0~L
\end{enumerate}
\end{minipage}\\[0.2cm]
\hline
\end{tabular}

\vspace{0.2cm}
\noindent\begin{tabular}{|p{0.9\linewidth}|}
\hline
\textbf{(35)} A 10.0~L rigid container holds 0.50~mol of oxygen gas at 1.5~atm. If the amount of gas doubles while temperature stays constant, what will the final pressure be?\\[0.2cm]
\begin{minipage}[t]{0.85\linewidth}
\begin{enumerate}[label=\alph*.]
\item 0.75~atm
\item 1.5~atm
\item 3.0~atm
\item 6.0~atm
\end{enumerate}
\end{minipage}\\[0.2cm]
\hline
\end{tabular}

\vspace{0.2cm}
\noindent\begin{tabular}{|p{0.9\linewidth}|}
\hline
\textbf{(36)} A weather balloon has a volume of 50.0~L at 250.~K. If it rises and the temperature drops to 200.~K while pressure stays constant, what is the new volume?\\[0.2cm]
\begin{minipage}[t]{0.85\linewidth}
\begin{enumerate}[label=\alph*.]
\item 25.0~L
\item 40.0~L
\item 50.0~L
\item 62.5~L
\end{enumerate}
\end{minipage}\\[0.2cm]
\hline
\end{tabular}
\newpage

\vspace{0.2cm}
\noindent\begin{tabular}{|p{0.9\linewidth}|}
\hline
\textbf{(37)} A tire is inflated with 3.0~mol of air to a volume of 18.0~L. If an additional 1.0~mol of air is added without changing temperature or pressure, what is the final volume?\\[0.2cm]
\begin{minipage}[t]{0.85\linewidth}
\begin{enumerate}[label=\alph*.]
\item 13.5~L
\item 18.0~L
\item 24.0~L
\item 30.0~L
\end{enumerate}
\end{minipage}\\[0.2cm]
\hline
\end{tabular}

\noindent\begin{tabular}{|p{0.9\linewidth}|}
\hline
\textbf{(38)} A balloon contains 1.00~mol of helium and occupies 24.5~L at room conditions. If 2.00~mol of helium are added at the same temperature and pressure, what is the new volume of the balloon?\\[0.2cm]
\begin{minipage}[t]{0.85\linewidth}
\begin{enumerate}[label=\alph*.]
\item 12.3~L
\item 24.5~L
\item 36.8~L
\item 49.0~L
\end{enumerate}
\end{minipage}\\[0.2cm]
\hline
\end{tabular}

\vspace{0.2cm}
\noindent\begin{tabular}{|p{0.9\linewidth}|}
\hline
\textbf{(39)} Two similar compounds, CH$_3$CH$_2$CH$_2$OH and CH$_3$OCH$_2$CH$_3$, have different boiling points. The alcohol boils at a much higher temperature. Why?\\[0.2cm]
\begin{minipage}[t]{0.85\linewidth}
\begin{enumerate}[label=\alph*.]
\item The alcohol has hydrogen bonding, while the ether does not.
\item The ether has stronger dipole-dipole forces.
\item Both compounds have identical London dispersion forces.
\item The ether forms ionic interactions in liquid form.
\end{enumerate}
\end{minipage}\\[0.2cm]
\hline
\end{tabular}

\noindent\begin{tabular}{|p{0.9\linewidth}|}
\hline
\textbf{(40)} Maya accidentally spills water and ethanol on the same counter. She notices that the ethanol evaporates faster. Which statement best explains why?\\[0.2cm]
\begin{minipage}[t]{0.85\linewidth}
\begin{enumerate}[label=\alph*.]
\item Ethanol molecules have stronger hydrogen bonding than water. 
\item Water has stronger hydrogen bonding, which slows evaporation.
\item Ethanol is nonpolar, so it doesn't interact with itself.
\item Water molecules lack dipole-dipole forces.
\end{enumerate}
\end{minipage}\\[0.2cm]
\hline
\end{tabular}
\newpage

\vspace{0.2cm}
\noindent\begin{tabular}{|p{0.9\linewidth}|}
\hline
\textbf{(41)} Two rigid containers are connected by a closed valve. One container holds helium at 2.0~atm in 4.0~L, and the other holds neon at 4.0~atm in 2.0~L. When the valve is opened and gases mix, what is the total pressure in the combined 6.0~L system (assume constant temperature)?\\[0.2cm]
\begin{minipage}[t]{0.85\linewidth}
\begin{enumerate}[label=\alph*.]
\item 2.0~atm
\item 3.3~atm
\item 4.0~atm
\item 6.0~atm
\end{enumerate}
\end{minipage}\\[0.2cm]
\hline
\end{tabular}

\vspace{0.2cm}
\noindent\begin{tabular}{|p{0.9\linewidth}|}
\hline
\textbf{(42)} When cooking pasta, steam rises from the boiling water. Which intermolecular forces are overcome as the water vaporizes?\\[0.2cm]
\begin{minipage}[t]{0.85\linewidth}
\begin{enumerate}[label=\alph*.]
\item Hydrogen bonding between water molecules
\item Ionic attractions between H$^+$ and O$^{2-}$
\item Metallic bonding between hydrogen atoms
\item Covalent bonds within H$_2$O molecules
\end{enumerate}
\end{minipage}\\[0.2cm]
\hline
\end{tabular}

\noindent\begin{tabular}{|p{0.9\linewidth}|}
\hline
\textbf{(43)} A solution is prepared by dissolving 10.0~g of NaCl in 90.0~g of water. What is the mass percent of NaCl in the solution?\\[0.2cm]
\begin{minipage}[t]{0.85\linewidth}
\begin{enumerate}[label=\alph*.]
\item 10.0\%
\item 11.1\%
\item 9.00\%
\item 12.0\%
\end{enumerate}
\end{minipage}\\[0.2cm]
\hline
\end{tabular}

\vspace{0.2cm}
\noindent\begin{tabular}{|p{0.9\linewidth}|}
\hline
\textbf{(44)} A sample of gas occupies 8.0~L at 1.0~atm. The gas is compressed until its volume is 2.0~L. What is the new pressure, assuming temperature is constant?\\[0.2cm]
\begin{minipage}[t]{0.85\linewidth}
\begin{enumerate}[label=\alph*.]
\item 0.25~atm
\item 2.0~atm
\item 4.0~atm
\item 8.0~atm
\end{enumerate}
\end{minipage}\\[0.2cm]
\hline
\end{tabular}
\newpage

\vspace{0.2cm}
\noindent\begin{tabular}{|p{0.9\linewidth}|}
\hline
\textbf{(45)} A balloon containing 2.50~mol of nitrogen occupies 60.0~L. If another 1.50~mol of nitrogen are added at constant temperature and pressure, what is the final volume?\\[0.2cm]
\begin{minipage}[t]{0.85\linewidth}
\begin{enumerate}[label=\alph*.]
\item 36.0~L
\item 60.0~L
\item 96.0~L
\item 120.~L
\end{enumerate}
\end{minipage}\\[0.2cm]
\hline
\end{tabular}

\vspace{0.2cm}
\noindent\begin{tabular}{|p{0.9\linewidth}|}
\hline
\textbf{(46)} According to solubility rules, most nitrate (NO$_3^-$) salts are soluble in water.\\[0.2cm]
\textbf{True} \hspace{1cm} \textbf{False}\\[0.2cm]
\hline
\end{tabular}

\vspace{0.2cm}
\noindent\begin{tabular}{|p{0.9\linewidth}|}
\hline
\textbf{(47)} A lab technician finds that NaCl dissolves easily in water but not in hexane. What is the main reason for this?\\[0.2cm]
\begin{minipage}[t]{0.85\linewidth}
\begin{enumerate}[label=\alph*.]
\item Water's polarity allows it to interact strongly with ionic compounds.
\item Hexane forms hydrogen bonds with Na$^+$ ions.
\item NaCl forms London forces with water.
\item Hexane is polar and repels ions.
\end{enumerate}
\end{minipage}\\[0.2cm]
\hline
\end{tabular}

\vspace{0.2cm}
\noindent\begin{tabular}{|p{0.9\linewidth}|}
\hline
\textbf{(48)} A scientist collects hydrogen gas over water at 25.0~°C. The total pressure is 745~torr, and the vapor pressure of water at 25.0~°C is 24.0~torr. What is the partial pressure of hydrogen?\\[0.2cm]
\begin{minipage}[t]{0.85\linewidth}
\begin{enumerate}[label=\alph*.]
\item 721~torr
\item 769~torr
\item 24.0~torr
\item 745~torr
\end{enumerate}
\end{minipage}\\[0.2cm]
\hline
\end{tabular}
\newpage

\vspace{0.2cm}
\noindent\begin{tabular}{|p{0.9\linewidth}|}
\hline
\textbf{(49)} A gas syringe contains 50.0~mL of nitrogen at 750~torr. If the gas is compressed to 25.0~mL, what is the final pressure (in torr) if temperature remains constant?\\[0.2cm]
\begin{minipage}[t]{0.85\linewidth}
\begin{enumerate}[label=\alph*.]
\item 375~torr
\item 500~torr
\item 1000~torr
\item 1500~torr
\end{enumerate}
\end{minipage}\\[0.2cm]
\hline
\end{tabular}

\noindent\begin{tabular}{|p{0.9\linewidth}|}
\hline
\textbf{(50)} A balloon is filled with a gas mixture containing 0.6~mol O$_2$ and 0.4~mol N$_2$ at a total pressure of 1.0~atm. What is the partial pressure of O$_2$?\\[0.2cm]
\begin{minipage}[t]{0.85\linewidth}
\begin{enumerate}[label=\alph*.]
\item 0.4~atm
\item 0.6~atm
\item 1.0~atm
\item 1.5~atm
\end{enumerate}
\end{minipage}\\[0.2cm]
\hline
\end{tabular}

\vspace{0.2cm}
\noindent\begin{tabular}{|p{0.9\linewidth}|}
\hline
\textbf{(51)} Dry ice (solid CO$_2$) sublimes directly into a gas at room temperature. Which type of intermolecular force is responsible for holding its molecules together in the solid phase?\\[0.2cm]
\begin{minipage}[t]{0.85\linewidth}
\begin{enumerate}[label=\alph*.]
\item Ionic forces
\item Hydrogen bonding
\item London dispersion forces
\item Dipole-dipole interactions
\end{enumerate}
\end{minipage}\\[0.2cm]
\hline
\end{tabular}

\noindent\begin{tabular}{|p{0.9\linewidth}|}
\hline
\textbf{(52)} Liquid nitrogen boils at -196~°C, while water boils at 100~°C. What difference in intermolecular forces best explains this vast difference?\\[0.2cm]
\begin{minipage}[t]{0.85\linewidth}
\begin{enumerate}[label=\alph*.]
\item Nitrogen has hydrogen bonding; water has dispersion only.
\item Water forms hydrogen bonds; nitrogen has only London forces.
\item Both have identical intermolecular forces.
\item Nitrogen has dipole-dipole interactions stronger than water. 
\end{enumerate}
\end{minipage}\\[0.2cm]
\hline
\end{tabular}
\newpage

\vspace{0.2cm}
\noindent\begin{tabular}{|p{0.9\linewidth}|}
\hline
\textbf{(53)} A sample of hydrogen occupies 10.0~L when it contains 0.400~mol of gas. What would be the volume if the amount were increased to 1.00~mol under the same conditions?\\[0.2cm]
\begin{minipage}[t]{0.85\linewidth}
\begin{enumerate}[label=\alph*.]
\item 4.00~L
\item 10.0~L
\item 25.0~L
\item 40.0~L
\end{enumerate}
\end{minipage}\\[0.2cm]
\hline
\end{tabular}

\vspace{0.2cm}
\noindent\begin{tabular}{|p{0.9\linewidth}|}
\hline
\textbf{(54)} A sealed cylinder of gas has a volume of 8.0~L at 280~K. If it is heated to 350~K at constant pressure, what is the final volume?\\[0.2cm]
\begin{minipage}[t]{0.85\linewidth}
\begin{enumerate}[label=\alph*.]
\item 6.4~L
\item 8.0~L
\item 10.0~L
\item 12.0~L
\end{enumerate}
\end{minipage}\\[0.2cm]
\hline
\end{tabular}

\vspace{0.2cm}
\noindent\begin{tabular}{|p{0.9\linewidth}|}
\hline
\textbf{(55)} A rigid 2.0~L cylinder contains 3.0~mol of gas at 5.0~atm. Some gas leaks until only 1.5~mol remain, while temperature and volume remain unchanged. What is the new pressure?\\[0.2cm]
\begin{minipage}[t]{0.85\linewidth}
\begin{enumerate}[label=\alph*.]
\item 1.5~atm
\item 2.5~atm
\item 5.0~atm
\item 10.~atm
\end{enumerate}
\end{minipage}\\[0.2cm]
\hline
\end{tabular}

\vspace{0.2cm}
\noindent\begin{tabular}{|p{0.9\linewidth}|}
\hline
\textbf{(56)} A 15.0~L balloon contains 0.600~mol of helium. If the balloon is filled until it contains 1.50~mol of helium, what will the new volume be at the same temperature and pressure?\\[0.2cm]
\begin{minipage}[t]{0.85\linewidth}
\begin{enumerate}[label=\alph*.]
\item 6.00~L
\item 15.0~L
\item 25.0~L
\item 37.5~L
\end{enumerate}
\end{minipage}\\[0.2cm]
\hline
\end{tabular}
\newpage

\vspace{0.2cm}
\noindent\begin{tabular}{|p{0.9\linewidth}|}
\hline
\textbf{(57)} A solution contains 0.50~mol of NaCl in 250~mL of solution. What is the molarity?\\[0.2cm]
\begin{minipage}[t]{0.85\linewidth}
\begin{enumerate}[label=\alph*.]
\item 0.50~M
\item 1.0~M
\item 2.0~M
\item 0.25~M
\end{enumerate}
\end{minipage}\\[0.2cm]
\hline
\end{tabular}

\vspace{0.2cm}
\noindent\begin{tabular}{|p{0.9\linewidth}|}
\hline
\textbf{(58)} A 10.0~L container holds 0.500~mol of gas at constant temperature and pressure. What volume would 2.00~mol of the same gas occupy under identical conditions?\\[0.2cm]
\begin{minipage}[t]{0.85\linewidth}
\begin{enumerate}[label=\alph*.]
\item 2.50~L
\item 5.00~L
\item 10.0~L
\item 40.0~L
\end{enumerate}
\end{minipage}\\[0.2cm]
\hline
\end{tabular}

\vspace{0.2cm}
\noindent\begin{tabular}{|p{0.9\linewidth}|}
\hline
\textbf{(59)} A chemist dissolves 50~mL of ethanol in 150~mL of water. What is the percent by volume of ethanol?\\[0.2cm]
\begin{minipage}[t]{0.85\linewidth}
\begin{enumerate}[label=\alph*.]
\item 25\%
\item 33.3\%
\item 50\%
\item 66.7\%
\end{enumerate}
\end{minipage}\\[0.2cm]
\hline
\end{tabular}

\vspace{0.2cm}
\noindent\begin{tabular}{|p{0.9\linewidth}|}
\hline
\textbf{(60)} What is the mass percent of glucose in a solution prepared by dissolving 18~g of glucose in 82~g of water?\\[0.2cm]
\begin{minipage}[t]{0.85\linewidth}
\begin{enumerate}[label=\alph*.]
\item 18\%
\item 20\%
\item 22\%
\item 25\%
\end{enumerate}
\end{minipage}\\[0.2cm]
\hline
\end{tabular}

\newpage
{\large\textbf{Checking for Understanding (general conceptual questions)}}\\
\noindent \begin{tabular}{|p{0.9\linewidth}|}
\hline
\vspace{0.2cm}

\begin{enumerate}
    \item \textbf{Targeted focus through self-selection:} When students choose which questions to tackle from the worksheet, they naturally gravitate toward areas they feel less confident about, ensuring the group addresses gaps in understanding.
    \item \textbf{Peer modeling of problem-solving:} Seeing a peer work through a problem on the board exposes students to different ways of setting up and explaining a solution, clarifying concepts for observers.
    \item \textbf{Reflection on understanding:} Students must decide whether they can solve a question themselves or need help, prompting metacognition and highlighting what they truly understand versus what requires clarification.
    \item \textbf{Immediate correction of misconceptions:} Peers correcting or expanding an explanation in real time prevents misunderstandings from becoming ingrained.
    \item \textbf{Balancing independence and collaboration:} Students first work independently on a question, then collaborate during discussion, ensuring individual accountability while leveraging group problem-solving.
    \item \textbf{Distinguishing conceptual vs. procedural questions:} Students can classify questions by identifying key ideas versus stepwise calculations, determining whether a question emphasizes big-picture reasoning or step-by-step execution.
    \item \textbf{Connecting big concepts to details:} Rotating through different student-selected questions helps the group link overarching concepts to multiple specific examples, reinforcing understanding of the framework behind the problems.
    \item \textbf{Verbalizing reasoning:} Explaining steps out loud forces students to articulate assumptions, logic, and chemical principles, improving clarity and communication of ideas.
    \item \textbf{Learning from multiple approaches:} Different students may solve the same problem uniquely, exposing peers to alternative methods or perspectives and broadening their problem-solving toolkit.
    \item \textbf{Application of theory to practice:} Working collaboratively on worksheet problems allows students to connect abstract chemical principles to concrete calculations or predictions.
    \item \textbf{Instructor insight:} Questions students select or ask reveal where common difficulties lie, helping instructors or SI leaders provide targeted guidance.
    \item \textbf{Strengthening long-term understanding:} Actively solving, explaining, and correcting questions in real time consolidates memory and understanding more effectively than passive review.
\end{enumerate}
\hline
\end{tabular}

\newpage
\noindent
\newpage
\noindent
\section*{\underline{Activity 2}}
\begin{tabular}{|p{4.2cm}|p{9.8cm}|}
\hline
\textbf{Collaborative Learning Techniques} & Clusters \\
\hline
\textbf{Learning Strategy} & Escape Room \\
\hline
\textbf{Time} & 12:35---12:50 \\
\hline
\textbf{Materials} & A periodic table, the textbook, readiness to work as a team or share ideas \\
\hline
\textbf{Lecture/Textbook/Old Plans/Slide References} & Chapter 7: Gases, Liquids and Solids, Chapter 8: Solutions \\
\hline
\end{tabular}

\vspace{0.5cm}
\noindent
\begin{tabular}{|p{0.9\linewidth}|}
\hline
{\large\textbf{Definitions}}
\begin{enumerate}
    \item \textbf{Escape Room:} Students enter a themed, time-bound challenge where they must solve a series of puzzles, riddles, or tasks to “escape” the room. Each puzzle is designed to test understanding of specific course concepts, requiring critical thinking, teamwork, and creativity. Students can expect to encounter hidden clues, coded messages, or problem-solving stations that build on one another; progress usually depends on collaboration and discussion. The SI leader will circulate to provide subtle hints, clarify misconceptions, and ensure students stay on track without giving away answers. This activity promotes peer-to-peer learning, reinforces content in an engaging way, and encourages students to communicate their reasoning while practicing time management under pressure.
\end{enumerate} \\
\hline
\end{tabular}
\newpage

\vspace{0.5cm}
\noindent
{\large \textbf{Instructions (detailed activity breakdown below)}} \\

\begin{tabular}{|p{0.9\linewidth}|}
\hline
\begin{enumerate}
    \item \textbf{Objective:} Students will collaboratively solve chemistry problems covering \textbf{solutions, gas laws, KMT, and intermolecular forces}. Each correct response earns points toward an overall team score. The goal is to promote discussion, problem-solving, and application of multiple concepts in a dynamic setting

    \item \textbf{Group Size:} Pairs (2 students per team)

    \item \textbf{Setup (Teacher Preparation):}
    \begin{itemize}
        \item Arrange \textbf{5–6 stations}, each focused on one major topic area:
        \begin{itemize}
            \item \textbf{Station 1 – Charles' Law}
            \item \textbf{Station 2 – Dalton's Law}
            \item \textbf{Station 3 – Boyle's Law}
            \item \textbf{Station 4 – Avogadro's Law}
            \item \textbf{Station 5 - Ideal gas Law}
            \item \textbf{Station 6 - Concentration}
            \item \textbf{Station 7 - Solubility}
            \item \textbf{(Optional) Station 8 – Hint / Checkpoint Station}
        \end{itemize}
        \item At each station, place a set of laminated question cards (multiple choice and true/false). Each card includes:
        \begin{itemize}
            \item A unique question number and topic label.
            \item The possible answer choices (A–E or T/F).
            \item A mapping line showing which station each answer leads to (to simulate “pathways” through the escape room).
        \end{itemize}
        \item Provide each pair of students with a \textbf{team slip} (half-sheet) labeled with their team name and space to record up to 3 verified answers.
    \end{itemize}
\end{enumerate}
\\
\hline
\end{tabular}
\newpage

\vspace{0.5cm}
\noindent
{\large \textbf{Instructions (detailed activity breakdown below) conitnued...}} \\

\begin{tabular}{|p{0.9\linewidth}|}
\hline
\begin{enumerate}
    \item \textbf{Facilitation:}
    \begin{itemize}
        \item Pairs may begin at \textbf{any station of their choice}.
        \item Each pair selects one question card, discusses the problem, and writes their chosen answer \textbf{on the back of the card}.
        \item Students then show their work and reasoning to the instructor or SI leader for quick verification.
        \item If the answer is correct, they record that problem on their team slip (problem number and answer) and may proceed to any other station of their choice.
        \item Pairs are encouraged to move freely between stations, aiming to correctly solve at least \textbf{three problems} during the activity period.
        \item After completing at least three problems, teams submit their slips to the instructor for scoring.
    \end{itemize}

    \item \textbf{Scoring:}
    \begin{itemize}
        \item Each correctly solved problem = \textbf{1 point}.
        \item Teams with the highest total score win or “escape” first.
        \item Bonus points may be awarded for clear reasoning or teamwork.
    \end{itemize}

    \item \textbf{Rules:}
    \begin{itemize}
        \item Teams must work collaboratively—no switching partners mid-activity.
        \item All work must be shown before a response is accepted.
        \item Only verified answers written on the team slip count toward the final score.
        \item Communication between pairs is encouraged, but direct sharing of answers is not permitted.
    \end{itemize}

    \item \textbf{Wrap-Up:}
    \begin{itemize}
        \item After all slips are collected, review several challenging questions as a class.
        \item Discuss common reasoning errors and strategies used by top-scoring pairs.
        \item Reinforce conceptual links between stations (e.g., gas behavior and intermolecular forces).
    \end{itemize}
\end{enumerate}
\\
\hline
\end{tabular}



\newpage
\noindent
{\large\textbf{Checking for Understanding (questions \& answers)}}\\[0.2em]
\noindent
\begin{tabular}{|p{0.9\linewidth}|}
\hline
\textbf{Questions:} \\
\begin{enumerate}
    \item \textbf{Targeted clarification:} Allowing students to bring questions from the lecture lets them focus on concepts they found most confusing.
    \item \textbf{Peer modeling of explanations:} Observing how classmates summarize or interpret lecture content helps students see alternative ways to understand the material.
    \item \textbf{Reflection on comprehension:} Discussing lecture points collaboratively encourages students to distinguish what they grasp from what needs further clarification.
    \item \textbf{Immediate correction of misconceptions:} Peers can correct or expand on explanations in real time, preventing misunderstandings from solidifying.
    \item \textbf{Balancing independent thinking and discussion:} Students first reflect individually, then contribute to group discussion, promoting both accountability and collaborative learning.
    \item \textbf{Identifying conceptual vs. procedural points:} Students can highlight which parts of the lecture are big-picture ideas versus detailed steps or examples.
    \item \textbf{Connecting concepts across topics:} Discussing multiple lecture points together helps students link overarching principles to specific details.
    \item \textbf{Verbalizing reasoning:} Explaining ideas out loud strengthens students’ ability to communicate and reason about the material clearly.
    \item \textbf{Learning from diverse perspectives:} Hearing multiple interpretations or summaries of the same lecture point exposes students to alternative approaches.
    \item \textbf{Application of theory:} Collaborative discussion encourages students to connect lecture content to examples, problems, or real-world applications.
    \item \textbf{Instructor insight:} Observing which lecture points students question most helps the instructor identify common difficulties and adjust instruction.
    \item \textbf{Reinforcing long-term understanding:} Actively reviewing and discussing lecture material together solidifies comprehension better than passive note review.
\end{enumerate} \\
\hline  
\end{tabular}

\newpage
\noindent
\section*{\underline{Closer}}
\begin{tabular}{|p{4.2cm}|p{9.8cm}|}
\hline
\textbf{Collaborative Learning Techniques} & Group Survey\\
\hline
\textbf{Learning Strategy} & Brain dump\\
\hline
\textbf{Time} & 12:50---12:55\\
\hline
\textbf{Materials} & Index cards, writing utensils\\
\hline
\textbf{Lecture/Textbook/Old Plans/Slide References} & Chapters 1-8 \\
\hline
\end{tabular}

\vspace{0.5cm}
\noindent
\begin{tabular}{|p{0.9\linewidth}|}
\hline
{\large\textbf{Definitions}}
\begin{enumerate}
   \item\textbf{Group survey:} Each group member is surveyed to discover their position on an issue, problem or topic. This process ensures that each member of the group is allowed to offer or state their point of view.
   \item\textbf{Brain dump:} This is a great closer for students to sum up their knowledge on the topic(s) from the session. Have the students take out a blank piece of paper /Google Doc and write/type("dump") everything they know relevant from the discussion that day. The SI Leader can look at their work to assess the students' understanding of the course material and have them compare it to each other. It is important to at least discuss with the students so the SI Leader knows the students' comfort levels with the material. 
\end{enumerate}
\\ \hline
\end{tabular}

\newpage
\noindent
{\large\textbf{Instructions (please provide a\hl{DETAILED} activity breakdown below (and/or attached}}

\begin{tabular}{|p{0.9\linewidth}|}
\hline
\begin{enumerate}[itemsep=0.3em]
   \item Students are put into small groups and are asked to write down their confidence on a single notecard without names.
   \item The cards are evidence of the group's confidence in the content.
\end{enumerate}
\\ \hline
\end{tabular}

\vspace{0.5cm}
\noindent
{\large\textbf{Content (fully worked out problems \& answers)}

\begin{tabular}{|p{0.9\linewidth}|}
\hline
\begin{itemize}
   \item No new worked problems; brain dump.
\end{itemize}
\\ \hline
\end{tabular}

\vspace{0.5cm}
\noindent
{\large\textbf{Checking for Understanding (questions \& answers) continued...}

\begin{tabular}{|p{0.9\linewidth}|}
\hline
\begin{itemize}
   \item Review confidence levels
   \item Interpret the variety of confidence and their distribution in groups.
   \item Get feedback for what they like and don't like
   \item Check in if the time is suitable for the students
\end{itemize}
\\ \hline
\end{tabular}
\newpage
\section*{Plan Checklist}

\begin{tabular}{|p{0.9\linewidth}|}
\hline
\begin{itemize}
\item[$\Box$] Variety of Collaborative Learning Techniques and Learning Strategies
\item[$\Box$] Chapter/slide\# where information is coming from
\item[$\Box$] Included timestamps (and buffer time as needed)
\item[$\Box$] Details on how leader will break up students
\item[$\Box$] Checking for Understanding questions
\item[$\Box$] Fully worked out problems and examples for each activity
\item[$\Box$] Necessary links required for online implementation (if applicable)
\item[$\Box$] Source material correctly cited (if applicable)
\item[$\Box$] Plan is uploaded to Teams
\end{itemize}
\\ \hline
\end{tabular}
\newpage

\end{document}